% Options for packages loaded elsewhere
% Options for packages loaded elsewhere
\PassOptionsToPackage{unicode}{hyperref}
\PassOptionsToPackage{hyphens}{url}
\PassOptionsToPackage{dvipsnames,svgnames,x11names}{xcolor}
%
\documentclass[
  letterpaper,
  DIV=11,
  numbers=noendperiod,
  oneside]{scrartcl}
\usepackage{xcolor}
\usepackage[left=1in,marginparwidth=2.0666666666667in,textwidth=4.1333333333333in,marginparsep=0.3in]{geometry}
\usepackage{amsmath,amssymb}
\setcounter{secnumdepth}{-\maxdimen} % remove section numbering
\usepackage{iftex}
\ifPDFTeX
  \usepackage[T1]{fontenc}
  \usepackage[utf8]{inputenc}
  \usepackage{textcomp} % provide euro and other symbols
\else % if luatex or xetex
  \usepackage{unicode-math} % this also loads fontspec
  \defaultfontfeatures{Scale=MatchLowercase}
  \defaultfontfeatures[\rmfamily]{Ligatures=TeX,Scale=1}
\fi
\usepackage{lmodern}
\ifPDFTeX\else
  % xetex/luatex font selection
\fi
% Use upquote if available, for straight quotes in verbatim environments
\IfFileExists{upquote.sty}{\usepackage{upquote}}{}
\IfFileExists{microtype.sty}{% use microtype if available
  \usepackage[]{microtype}
  \UseMicrotypeSet[protrusion]{basicmath} % disable protrusion for tt fonts
}{}
\makeatletter
\@ifundefined{KOMAClassName}{% if non-KOMA class
  \IfFileExists{parskip.sty}{%
    \usepackage{parskip}
  }{% else
    \setlength{\parindent}{0pt}
    \setlength{\parskip}{6pt plus 2pt minus 1pt}}
}{% if KOMA class
  \KOMAoptions{parskip=half}}
\makeatother
% Make \paragraph and \subparagraph free-standing
\makeatletter
\ifx\paragraph\undefined\else
  \let\oldparagraph\paragraph
  \renewcommand{\paragraph}{
    \@ifstar
      \xxxParagraphStar
      \xxxParagraphNoStar
  }
  \newcommand{\xxxParagraphStar}[1]{\oldparagraph*{#1}\mbox{}}
  \newcommand{\xxxParagraphNoStar}[1]{\oldparagraph{#1}\mbox{}}
\fi
\ifx\subparagraph\undefined\else
  \let\oldsubparagraph\subparagraph
  \renewcommand{\subparagraph}{
    \@ifstar
      \xxxSubParagraphStar
      \xxxSubParagraphNoStar
  }
  \newcommand{\xxxSubParagraphStar}[1]{\oldsubparagraph*{#1}\mbox{}}
  \newcommand{\xxxSubParagraphNoStar}[1]{\oldsubparagraph{#1}\mbox{}}
\fi
\makeatother


\usepackage{longtable,booktabs,array}
\newcounter{none} % for unnumbered tables
\usepackage{calc} % for calculating minipage widths
% Correct order of tables after \paragraph or \subparagraph
\usepackage{etoolbox}
\makeatletter
\patchcmd\longtable{\par}{\if@noskipsec\mbox{}\fi\par}{}{}
\makeatother
% Allow footnotes in longtable head/foot
\IfFileExists{footnotehyper.sty}{\usepackage{footnotehyper}}{\usepackage{footnote}}
\makesavenoteenv{longtable}
\usepackage{graphicx}
\makeatletter
\newsavebox\pandoc@box
\newcommand*\pandocbounded[1]{% scales image to fit in text height/width
  \sbox\pandoc@box{#1}%
  \Gscale@div\@tempa{\textheight}{\dimexpr\ht\pandoc@box+\dp\pandoc@box\relax}%
  \Gscale@div\@tempb{\linewidth}{\wd\pandoc@box}%
  \ifdim\@tempb\p@<\@tempa\p@\let\@tempa\@tempb\fi% select the smaller of both
  \ifdim\@tempa\p@<\p@\scalebox{\@tempa}{\usebox\pandoc@box}%
  \else\usebox{\pandoc@box}%
  \fi%
}
% Set default figure placement to htbp
\def\fps@figure{htbp}
\makeatother





\setlength{\emergencystretch}{3em} % prevent overfull lines

\providecommand{\tightlist}{%
  \setlength{\itemsep}{0pt}\setlength{\parskip}{0pt}}



 


\KOMAoption{captions}{tableheading}
\makeatletter
\@ifpackageloaded{caption}{}{\usepackage{caption}}
\AtBeginDocument{%
\ifdefined\contentsname
  \renewcommand*\contentsname{Table of contents}
\else
  \newcommand\contentsname{Table of contents}
\fi
\ifdefined\listfigurename
  \renewcommand*\listfigurename{List of Figures}
\else
  \newcommand\listfigurename{List of Figures}
\fi
\ifdefined\listtablename
  \renewcommand*\listtablename{List of Tables}
\else
  \newcommand\listtablename{List of Tables}
\fi
\ifdefined\figurename
  \renewcommand*\figurename{Figure}
\else
  \newcommand\figurename{Figure}
\fi
\ifdefined\tablename
  \renewcommand*\tablename{Table}
\else
  \newcommand\tablename{Table}
\fi
}
\@ifpackageloaded{float}{}{\usepackage{float}}
\floatstyle{ruled}
\@ifundefined{c@chapter}{\newfloat{codelisting}{h}{lop}}{\newfloat{codelisting}{h}{lop}[chapter]}
\floatname{codelisting}{Listing}
\newcommand*\listoflistings{\listof{codelisting}{List of Listings}}
\makeatother
\makeatletter
\makeatother
\makeatletter
\@ifpackageloaded{caption}{}{\usepackage{caption}}
\@ifpackageloaded{subcaption}{}{\usepackage{subcaption}}
\makeatother
\makeatletter
\@ifpackageloaded{sidenotes}{}{\usepackage{sidenotes}}
\@ifpackageloaded{marginnote}{}{\usepackage{marginnote}}
\makeatother
\usepackage{bookmark}
\IfFileExists{xurl.sty}{\usepackage{xurl}}{} % add URL line breaks if available
\urlstyle{same}
\hypersetup{
  pdftitle={Chapter 2: Discipline},
  colorlinks=true,
  linkcolor={blue},
  filecolor={Maroon},
  citecolor={Blue},
  urlcolor={Blue},
  pdfcreator={LaTeX via pandoc}}


\title{Chapter 2: Discipline}
\usepackage{etoolbox}
\makeatletter
\providecommand{\subtitle}[1]{% add subtitle to \maketitle
  \apptocmd{\@title}{\par {\large #1 \par}}{}{}
}
\makeatother
\subtitle{Sociology, Science, and Core Concepts?}
\author{}
\date{}
\begin{document}
\maketitle


\section{Introduction}\label{introduction}

Here is one (very negative) view about the current state of Sociology:

\begin{quote}
According to Ray Rodrigues, chancellor of the state university system,
Florida lawmakers began eyeing general education as a target of reform
after complaints from parents. Their children would come home for
Thanksgiving break during their first semester in college and would ``be
communicating cultural Marxism''\ldots Under SB 266 {[}a Florida state
law intended to eliminate ``woke'' ideology{]}, courses that count
toward general-education credit should ``promote and preserve the
constitutional republic'' through ``traditional, historically accurate''
coursework. The law also says that ``core courses'' may not ``distort
significant historical events,'' ``include a curriculum that teaches
identity politics,'' or be based on theories that systemic racism,
sexism, oppression, and privilege are inherent in American
institutions.\sidenote{\footnotesize \href{https://www-chronicle-com.proxy.library.stonybrook.edu/article/the-curricular-cull}{``Is
  Sociology Being Saved --- or Debased --- in Florida?''}
  \emph{Chronicle of Higher Education}, March 16, 2026.}
\end{quote}

There is much to unpack in this statement. For now, we will leave aside
the meaning of ``cultural Marxism.''\sidenote{\footnotesize I seriously doubt that
  Rodrigues and his allies are reading Gramsci, Lukacs, or Benjamin, and
  instead, are simply using ``cultural Marxism'' as a catch-all for
  ideologies they don't like. See A.J.A. Woods, \emph{The Cultural
  Marxism Conspiracy}, Verso, 2026.} We will also only briefly note that
colleges explicitly advertise themselves as places to ``push boundaries,
challenge assumptions and redefine what's possible,''\sidenote{\footnotesize \href{https://www.stonybrook.edu/}{Stony
  Brook's Webpage}, accessed June 27th, 2026.} so one would
\emph{expect} young people to \emph{want} to discuss new and interesting
ideas with their family. We will also only glancingly note that
``promot{[}ing{]} and preserv{[}ing{]} the constitutional republic'' is
\emph{itself} an ideological goal, and that pursuing it explicitly could
and would bias our understanding of society just like an explicitly
``woke'' one would.

Instead, I want to take aim at what seems to be a foundational
assumption of those criticizing sociology: that there is a way to pursue
a completely unbiased, objective, neutral study of society. The authors
of the Flordia state law imply, by claiming to cleanse University
curriculum of ``distortions,'' that there \emph{is} an objective, true,
and indisputable nature of society and history that is like the things
we know about the physical history of the universe, biological dynamics,
or the engineering of the material environment.

As you will see in this chapter, at a very basic level sociology (and
the rest of the social sciences) do not work this way. This is because
they take as their central object \emph{people in society}, and since
people \emph{react} to their understandings about the world and the
meaning they take from it, they need to be studied in a different way
than rocks, planets, and pathogens. However, this does \emph{not} mean
that the social sciences and particularly sociology are therefore simply
a realm where we assert opinions and call them facts; instead, it means
that we have a wide inventory of careful, systematic ways of
understanding the world, and our job here in part is to introduce them.

\section{What Is A Science?}\label{what-is-a-science}

To answer the question ``what kind of science is sociology?'', we need
to answer a prior question: just what is a \textbf{science}? Science is
fundamentally a systematic investigation of what the world is like, and
its organization strives to achieves two major things.

\texttt{no\ supernatural\ expalanation}

\texttt{no\ settling\ of\ theoretical\ or\ empirical\ disputes\ by\ appeals\ to\ only\ power\ or\ authority}

\subsection{What Is A Theory?}\label{what-is-a-theory}

\texttt{A\ theory\ is\ a\ statement\ about\ how\ at\ least\ two\ different\ elements\ of\ the\ world\ relate\ to\ one\ another.}

\texttt{They\ can\ vary\ from\ descriptive\ to\ highly\ formal\ statements.}

\subsection{What Is A Hypothesis?}\label{what-is-a-hypothesis}

\texttt{Scientists\ generally\ think\ that\ theories\ are\ best\ when\ they\ are\ parsimonious\ (containing\ a\ few\ simple\ elements)\ and\ general\ (able\ to\ explain\ at\ least\ a\ whole\ class\ of\ cases).\ One\ way\ to\ extend\ and\ deepen\ theories\ is\ to\ ask\ what\ they\ *predict*\ about\ the\ world.\ \ If\ that\ prediction\ comes\ to\ pass,\ we\ typically\ call\ the\ underlying\ theory\ *corroborated*,\ and\ when\ it\ does\ not,\ we\ typically\ call\ the\ prediction\ *falsified*.\ \ Collectively,\ this\ process\ of\ making\ guesses\ based\ on\ our\ theoretical\ knowledge\ and\ then\ trying\ them\ out\ in\ the\ world\ is\ what\ we\ call\ **hypothetical**\ reasoning.}

\subsection{What Is Evidence?}\label{what-is-evidence}

\subsection{What Is A Method?}\label{what-is-a-method}

\subsection{Some Tensions In A Pluralistic
Science}\label{some-tensions-in-a-pluralistic-science}

\section{Key Concepts in Contemporary
Sociology}\label{key-concepts-in-contemporary-sociology}

\subsection{Social Construction}\label{social-construction}

One of the simplest and most direct, but also most notorious, concepts
in sociology is social construction. If you have ever played any kind of
game, engaged in make-believe with a young child, read fiction, or
traveled in an unfamiliar culture, you have encountered social
constructions and been aware of doing so. But \emph{also}, if you ever
used a bank, been in love, celebrated a birthday, or spoken to another
person, you've also encountered a social construction, but perhaps with
less awareness.

\subsection{Social Order and Social
Change}\label{social-order-and-social-change}

\texttt{The\ two\ classic,\ "big"\ statements\ of\ this\ are\ Durkheim\textquotesingle{}s\ dissertation\ book\ and\ Tonnies.}

\subsection{Agency and Structure}\label{agency-and-structure}

\subsection{Micro, Macro, Meso}\label{micro-macro-meso}

\section{Key Theories in Contemporary
Sociology}\label{key-theories-in-contemporary-sociology}

\subsection{Structural/Functionalism}\label{structuralfunctionalism}

\subsection{Conflict Theory/Marxism}\label{conflict-theorymarxism}

\subsection{Symbolic Interactionism}\label{symbolic-interactionism}

\subsection{Practice Theory}\label{practice-theory}

\subsection{Challenges to the Current Theoretical
Moment}\label{challenges-to-the-current-theoretical-moment}




\end{document}
